\documentclass[11pt]{article}
\usepackage[margin=1in]{geometry}
\usepackage{amsmath,amssymb,amsthm,mathtools}
\usepackage[T1]{fontenc}
\usepackage[utf8]{inputenc}
\usepackage{enumitem}
\usepackage{booktabs,longtable,array}
\usepackage{hyperref}
\usepackage{xcolor}
\usepackage{microtype}
\hypersetup{colorlinks=true, linkcolor=blue!50!black, citecolor=blue!50!black, urlcolor=blue!50!black}

\newtheorem{theorem}{Theorem}[section]
\newtheorem{proposition}[theorem]{Proposition}
\newtheorem{lemma}[theorem]{Lemma}
\newtheorem{corollary}[theorem]{Corollary}
\newtheorem{claim}[theorem]{Claim}
\theoremstyle{definition}
\newtheorem{definition}[theorem]{Definition}
\newtheorem{assumption}[theorem]{Assumption}
\theoremstyle{remark}
\newtheorem{remark}[theorem]{Remark}

\newcommand{\KL}{\mathrm{KL}}
\newcommand{\Ent}{\mathrm{Ent}}
\newcommand{\Prob}{\mathsf{P}}
\newcommand{\E}{\mathbb{E}}
\newcommand{\R}{\mathbb{R}}
\newcommand{\N}{\mathbb{N}}
\newcommand{\Pcal}{\mathcal{P}}
\newcommand{\Mcal}{\mathcal{M}}
\newcommand{\Qcal}{\mathcal{Q}}
\newcommand{\Zcal}{\mathcal{Z}}
\newcommand{\Fcal}{\mathcal{F}}
\newcommand{\Gcal}{\mathcal{G}}
\newcommand{\W}{\mathcal{W}}
\newcommand{\one}{\mathbf{1}}
\newcommand{\dd}{\,d}
\newcommand{\supp}{\operatorname{supp}}
\newcommand{\argmin}{\operatorname*{arg\,min}}
\newcommand{\argmax}{\operatorname*{arg\,max}}
\newcommand{\esssup}{\operatorname*{ess\,sup}}
\newcommand{\ri}{\operatorname{ri}}
\newcommand{\dom}{\operatorname{dom}}
\newcommand{\Law}{\operatorname{Law}}
\newcommand{\Var}{\operatorname{Var}}
\newcommand{\Cov}{\operatorname{Cov}}
\newcommand{\diag}{\operatorname{diag}}
\newcommand{\Id}{\mathrm{Id}}
\newcommand{\Lip}{\operatorname{Lip}}
\newcommand{\osc}{\operatorname{osc}}
\newcommand{\diam}{\operatorname{diam}}
\newcommand{\proj}{\operatorname{proj}}
\newcommand{\e}{\mathrm{e}}

\newenvironment{argument}{\paragraph{Argument.}\begin{enumerate}[leftmargin=2em]}{\end{enumerate}}
\newenvironment{proofoutline}{\paragraph{Proof architecture.}\begin{enumerate}[leftmargin=2em]}{\end{enumerate}}


% Source-faithfulness apparatus used by the paper reconstructions.
\newcommand{\sourceanchor}[1]{\par\smallskip\noindent\textbf{Source anchor.} #1\par\smallskip}
\newcommand{\sourcecomment}[1]{\begin{remark}#1\end{remark}}
\newenvironment{sourceguide}
  {\begin{description}[style=nextline,leftmargin=3.2cm,labelwidth=3cm]}
  {\end{description}}

% Backward-compatible environments used by some of the older reconstructions.
\newenvironment{distilled}[1][]{\begin{theorem}[#1]}{\end{theorem}}
\newenvironment{distillation}{\begin{enumerate}[leftmargin=2em]}{\end{enumerate}}

\newenvironment{commentarynotes}{\begin{remark}\begin{enumerate}[leftmargin=2em]}{\end{enumerate}\end{remark}}

\title{Wasserstein Distributionally Robust Optimization on a Polish Space}
\author{}
\date{Gao and Kleywegt (2023): Theorem 1, Corollaries 1--2, and the constructive proof}
\begin{document}
\maketitle

\section*{Source map and scope}
\begin{description}[style=nextline,leftmargin=3.2cm,labelwidth=3cm]
\item[Primary source] \texttt{\detokenize{gao-kleywegt-2023_DRO-with-wasserstein-distance__arXiv-1604.02199.pdf}}
\item[Coverage] Section 3: Proposition 1, Definitions 3--4, Lemmas 2--7, Proposition 2, Theorem 1, Corollaries 1--2, and the proof architecture for existence and structure of worst-case laws.
\item[Source anchors]
\begin{itemize}[leftmargin=2em]
  \item equations (Dual), (3)--(8), (19)--(27);
  \item Proposition 1 and Proposition 2;
  \item Lemmas 2--7;
  \item Theorem 1 and Corollaries 1--2.
\end{itemize}
\end{description}

This reconstruction follows the source order and proof technique.  It retains the source sign convention
\(\Phi(\lambda,\zeta)=\inf_\xi\{\lambda d^p(\xi,\zeta)-\Psi(\xi)\}\), so the robust envelope is
\(-\Phi\).  Measurable-selection and normal-integrand facts used by the paper are stated at the points
where they enter; they are not replaced by a different duality theorem.

\section{Primal and scalar dual problems}
Let \((\Xi,d)\) be Polish, \(p\in[1,\infty)\), \(\nu\in\mathcal P_p(\Xi)\), \(\theta>0\), and
\(\Psi\in L^1(\nu)\).  The distributionally robust value is
\[
 v_P=\sup\left\{\int_\Xi\Psi\,d\mu: \mu\in\mathcal P(\Xi),\ W_p(\mu,\nu)\le\theta\right\}.
\]
The one-dimensional dual proposed in the paper is
\begin{equation}\label{eq:gk-dual}
 v_D=\inf_{\lambda\ge0}h(\lambda),\qquad
 h(\lambda)=\lambda\theta^p-\int_\Xi\Phi(\lambda,\zeta)\,\nu(d\zeta),
\end{equation}
where
\begin{equation}\label{eq:gk-Phi}
 \Phi(\lambda,\zeta)=\inf_{\xi\in\Xi}
 \{\lambda d^p(\xi,\zeta)-\Psi(\xi)\}.
\end{equation}

\subsection{Weak duality: Proposition 1}
For every \(\lambda\ge0\), \(\xi,\zeta\in\Xi\),
\[
 \Psi(\xi)\le \lambda d^p(\xi,\zeta)-\Phi(\lambda,\zeta).
\]
If \(\pi\in\Pi(\mu,\nu)\) has \(\int d^p\,d\pi\le\theta^p\), integration gives
\[
 \int\Psi\,d\mu
 \le \lambda\theta^p-\int\Phi(\lambda,\zeta)\,\nu(d\zeta)=h(\lambda).
\]
Taking the primal supremum and then the infimum over \(\lambda\) yields
\(v_P\le v_D\).

\section{Growth rate and the geometry of near minimizers}
For \(\Phi(\lambda,\zeta)>-\infty\), define the limiting extreme costs among
\(\varepsilon\)-minimizers:
\begin{align}
 D^+(\lambda,\zeta)
 &=\lim_{\varepsilon\downarrow0}
   \sup\{d^p(\xi,\zeta):
     \lambda d^p(\xi,\zeta)-\Psi(\xi)
       \le\Phi(\lambda,\zeta)+\varepsilon\},\label{eq:gk-Dplus}\\
 D^-(\lambda,\zeta)
 &=\lim_{\varepsilon\downarrow0}
   \inf\{d^p(\xi,\zeta):
     \lambda d^p(\xi,\zeta)-\Psi(\xi)
       \le\Phi(\lambda,\zeta)+\varepsilon\}.\label{eq:gk-Dminus}
\end{align}
When the pointwise infimum is attained, \(D_0^+\) and \(D_0^-\) denote the corresponding
largest and smallest exact minimizing costs.  The source warns that \(D_0^\pm\) need not equal
\(D^\pm\): minimizing sequences may escape even when exact minimizers exist.

The growth threshold is
\begin{equation}\label{eq:gk-kappa}
 \kappa=\inf\left\{\lambda\ge0:
       \int\Phi(\lambda,\zeta)\,\nu(d\zeta)>-\infty\right\}.
\end{equation}
For unbounded \(\Xi\), Lemma 2 identifies \(\kappa\) with the asymptotic positive
\(d^p\)-growth of \(\Psi\): for every base point \(\zeta\),
\begin{equation}\label{eq:gk-growth}
 \kappa=\limsup_{d^p(\xi,\zeta)\to\infty}
 \frac{\max\{0,\Psi(\xi)-\Psi(\zeta)\}}{d^p(\xi,\zeta)}.
\end{equation}
The base point does not matter because the triangle inequality controls the ratio of
\(d^p(\xi,\zeta)\) and \(d^p(\xi,\zeta')\) at infinity.  The same lemma proves that
\(\kappa<\infty\) exactly when there are \(L,M>0\) and \(\zeta_0\) such that
\[
 \Psi(\xi)-\Psi(\zeta_0)\le Ld^p(\xi,\zeta_0)+M,
 \qquad \xi\in\Xi.
\]
The forward implication follows by choosing any \(\lambda>\kappa\) with finite integrated
regularization and extracting a point \(\zeta_0\) for which \(\Phi(\lambda,\zeta_0)>-\infty\).
The reverse implication follows from the \(p\)-moment of \(\nu\) and
\(d^p(\xi,\zeta_0)\le 2^{p-1}(d^p(\xi,\zeta)+d^p(\zeta,\zeta_0))\).

\section{Measurability and selectors: Lemma 3}
The paper works with the completion \(\mathcal B_\nu(\Xi)\).  It proves that
\(\Phi(\lambda,\cdot)\), \(D^\pm(\lambda,\cdot)\), and, when defined,
\(D_0^\pm(\lambda,\cdot)\) are \(\nu\)-measurable.  For \(\varepsilon,\delta>0\), define
\begin{align*}
 F^+_{\varepsilon,\delta}(\lambda,\zeta)
 &=\{\xi: \lambda d^p(\xi,\zeta)-\Psi(\xi)
       \le\Phi(\lambda,\zeta)+\varepsilon,
       \ d^p(\xi,\zeta)\ge D^+(\lambda,\zeta)-\delta\},\\
 F^-_{\varepsilon,\delta}(\lambda,\zeta)
 &=\{\xi: \lambda d^p(\xi,\zeta)-\Psi(\xi)
       \le\Phi(\lambda,\zeta)+\varepsilon,
       \ d^p(\xi,\zeta)\le D^-(\lambda,\zeta)+\delta\}.
\end{align*}
Whenever these sets are nonempty almost surely, measurable-selection results give selectors
\(T^+_{\varepsilon,\delta}\) and \(T^-_{\varepsilon,\delta}\).  The same construction applies to exact
argmin sets and to the two escape-at-infinity sets needed in the proofs of Proposition 2 and Lemma 7.
This is the only point at which the proof needs Polish-space descriptive-set theory.

\section{The regularization operator: Lemma 4}
There is a full-measure set \(B\) such that, for \(\zeta\in B\):
\begin{enumerate}[leftmargin=2em,label=(\roman*)]
\item \(\Phi(\cdot,\zeta)\) is nondecreasing, concave, and upper semicontinuous;
\item \(\Phi(\lambda,\zeta)>-\infty\) for every \(\lambda>\kappa\);
\item increasing the multiplier shrinks the outer near-minimizer radius:
\[
 \lambda_2>\lambda_1>\kappa
 \quad\Longrightarrow\quad
 D^+(\lambda_2,\zeta)\le D^-(\lambda_1,\zeta)
 \le D^+(\lambda_1,\zeta);
\]
\item the minimizing radii satisfy the coercive estimate
\begin{equation}\label{eq:gk-radius-bound}
 (\lambda_2-\lambda_1)D^+(\lambda_2,\zeta)
 \le -\Psi(\zeta)-\Phi(\lambda_1,\zeta);
\end{equation}
\item the one-sided slopes of \(\Phi\) are bracketed by limiting minimizing costs:
\begin{align}
 D^-(\lambda,\zeta)
 &\le \partial_-\Phi(\lambda,\zeta)
 \le \lim_{\lambda_1\uparrow\lambda}D^+(\lambda_1,\zeta),\label{eq:gk-left-slope}\\
 \lim_{\lambda_2\downarrow\lambda}D^-(\lambda_2,\zeta)
 &\le \partial_+\Phi(\lambda,\zeta)
 \le D^+(\lambda,\zeta).\label{eq:gk-right-slope}
\end{align}
\end{enumerate}
Concavity is immediate because \(\Phi\) is an infimum of affine functions of \(\lambda\).
For the derivative inequalities, compare a near minimizer at one multiplier with the value at a
nearby multiplier, divide by the multiplier increment, and then let the near-optimality error vanish.

\section{The scalar objective: Lemma 5}
The dual objective \(h\) in \eqref{eq:gk-dual} is infinite on \([0,\kappa)\), finite on
\((\kappa,\infty)\), convex, lower semicontinuous, and coercive:
\[
 h(\lambda)\longrightarrow+\infty\qquad(\lambda\to\infty).
\]
The coercivity proof uses \(\Phi(\lambda,\zeta)\le -\Psi(\zeta)\), obtained by testing
\(\xi=\zeta\), so
\[
 h(\lambda)\ge \lambda\theta^p+\int\Psi\,d\nu.
\]
Consequently \(h\) has a minimizer \(\lambda^*\in[\kappa,\infty)\).

\section{The case \(\kappa=\infty\): Proposition 2}
If \(\kappa=\infty\), then \(h(\lambda)=+\infty\) for every \(\lambda\), hence \(v_D=+\infty\).
To prove \(v_P=+\infty\), fix a multiplier \(a>0\) and a payoff level \(b>0\).  The definition of
\(\kappa\) implies that on a positive-measure set of base points one can select \(T(\zeta)\) with
\[
 \Psi(T(\zeta))-\Psi(\zeta)>a d^p(T(\zeta),\zeta)+b.
\]
Move only a small fraction \(q\) of \(\nu\) along this selector and leave the remaining mass fixed:
\[
 \mu_q=qT_\#\nu+(1-q)\nu.
\]
Choosing \(q\) so that \(q\int d^p(T(\zeta),\zeta)d\nu\le\theta^p\) preserves feasibility, while the
payoff gain is at least \(qa\int d^p+qb\).  Taking \(a\) and \(b\) successively large makes the
primal value infinite.

\section{Interior dual minimizers: Lemma 6}
Assume \(\lambda^*>\kappa\).  The first-order conditions for convex \(h\) are
\[
 \partial_-h(\lambda^*)\le0\le\partial_+h(\lambda^*).
\]
The paper interchanges one-sided differentiation and integration.  For the right derivative, if
\(\lambda_n\downarrow\lambda^*\), the concavity difference quotients
\[
 f_n(\zeta)=\frac{\Phi(\lambda_n,\zeta)-\Phi(\lambda^*,\zeta)}
                  {\lambda_n-\lambda^*}
\]
are nonnegative and increase pointwise, so monotone convergence applies.  For the left derivative,
choose \(\lambda_1>\kappa\); the corresponding quotients are dominated by
\((|\Phi(\lambda_1,\zeta)|+|\Phi(\lambda^*,\zeta)|)/(\lambda^*-\lambda_1)\), which is integrable.
Thus
\begin{equation}\label{eq:gk-integrated-bracket}
 \int\partial_+\Phi(\lambda^*,\zeta)d\nu
 \le \theta^p
 \le \int\partial_-\Phi(\lambda^*,\zeta)d\nu.
\end{equation}
Using \eqref{eq:gk-left-slope}--\eqref{eq:gk-right-slope}, choose
\(\lambda_1^\varepsilon\uparrow\lambda^*\) and
\(\lambda_2^\varepsilon\downarrow\lambda^*\) so that measurable near-minimizer selectors
\(T_1^\varepsilon,T_2^\varepsilon\) have costs
\[
 A_1^\varepsilon=\int d^p(T_1^\varepsilon(\zeta),\zeta)d\nu
 \quad\text{and}\quad
 A_2^\varepsilon=\int d^p(T_2^\varepsilon(\zeta),\zeta)d\nu
\]
with \(A_1^\varepsilon\) asymptotically above \(\theta^p\) and
\(A_2^\varepsilon\) asymptotically below \(\theta^p\).  Introduce the identity selector with zero
cost.  There are weights \(p_1^\varepsilon,p_2^\varepsilon,p_3^\varepsilon\ge0\), summing to one,
such that
\[
 p_1^\varepsilon A_1^\varepsilon+p_2^\varepsilon A_2^\varepsilon\le\theta^p
\]
and the weighted cost approaches \(\theta^p\) whenever \(\lambda^*>0\).  Define
\begin{equation}\label{eq:gk-mueps-interior}
 \mu^\varepsilon=p_1^\varepsilon(T_1^\varepsilon)_\#\nu
                  +p_2^\varepsilon(T_2^\varepsilon)_\#\nu
                  +p_3^\varepsilon\nu.
\end{equation}
The explicit coupling that chooses the corresponding map conditional on \(\zeta\) proves
\(W_p^p(\mu^\varepsilon,\nu)\le\theta^p\).  The near-minimizer inequalities give
\[
 \Psi(T_i^\varepsilon(\zeta))
 \ge \lambda_i^\varepsilon d^p(T_i^\varepsilon(\zeta),\zeta)
      -\Phi(\lambda_i^\varepsilon,\zeta)-\varepsilon.
\]
After weighting, integrating, and using continuity of the integrated regularization at
\(\lambda^*\), one obtains the source estimate
\begin{equation}\label{eq:gk-interior-eps}
 \int\Psi\,d\mu^\varepsilon
 \ge \lambda^*\theta^p-
       \int\Phi(\lambda^*,\zeta)d\nu-\varepsilon.
\end{equation}

\section{Boundary dual minimizers: Lemma 7}
Assume \(\lambda^*=\kappa<\infty\).  The difficulty is that a minimizing sequence for
\(\Phi(\kappa,\zeta)\) may escape to infinity.  The source splits into two subcases.

\subsection{The zero-growth boundary \(\kappa=0\)}
Choose \(\lambda_2^\varepsilon\downarrow0\), a small near-minimizer tolerance, and a low-cost
selector for \(\lambda_2^\varepsilon\).  Mix its push-forward with the identity law so that the
transport budget is respected.  Since the identity component contributes \(\int\Psi d\nu\) and the
transported component satisfies the regularization inequality, the resulting primal value tends to
\(-\int\Phi(0,\zeta)d\nu=h(0)\).

\subsection{The positive-growth boundary \(\kappa>0\)}
Choose \(\kappa_0<\kappa\).  The growth-rate characterization and Lemma 3(v) provide an escape
selector \(T^R(\kappa_0,\zeta)\) satisfying
\[
 \Psi(T^R(\kappa_0,\zeta))-\Psi(\zeta)
 \ge \kappa_0d^p(T^R(\kappa_0,\zeta),\zeta)
\]
and whose cost is at least a prescribed radius \(R\).  A second selector is a near minimizer for a
multiplier \(\lambda>\kappa\).  The paper chooses a mixing coefficient
\(q_R(\kappa_0,\lambda)\) so that the weighted expected cost is exactly \(\theta^p\).  The large
cost of the escape selector makes its weight small, while its payoff-to-cost ratio approaches
\(\kappa\).  Letting
\[
 \kappa_0\uparrow\kappa,
 \qquad \lambda\downarrow\kappa,
 \qquad R\to\infty
\]
produces feasible measures \(\mu^\varepsilon\) satisfying
\begin{equation}\label{eq:gk-boundary-eps}
 \int\Psi\,d\mu^\varepsilon
 \ge \kappa\theta^p-
       \int\Phi(\kappa,\zeta)d\nu-5\varepsilon.
\end{equation}
The constants in the source are bookkeeping artifacts of the simultaneous choices; only their
vanishing matters.

\section{Theorem 1: strong duality}
If \(\kappa<\infty\), Lemma 5 supplies a dual minimizer.  Lemma 6 applies when
\(\lambda^*>\kappa\), and Lemma 7 applies at the boundary.  In either case, for every
\(\varepsilon>0\) there is a feasible \(\mu^\varepsilon\) with
\[
 \int\Psi\,d\mu^\varepsilon\ge v_D-C\varepsilon.
\]
Hence \(v_P\ge v_D\).  Together with Proposition 1,
\begin{theorem}[Gao--Kleywegt Theorem 1]
If \(\kappa<\infty\), then
\[
 v_P=v_D<\infty.
\]
The scalar dual problem admits a minimizer \(\lambda^*\in[\kappa,\infty)\).
\end{theorem}
Combined with Proposition 2, this describes both finite- and infinite-value cases.

\section{Corollary 1: exact worst-case laws}
Assume now that \(\Psi\) is upper semicontinuous and every bounded subset of \(\Xi\) is totally
bounded.  Then for \(\lambda>\kappa\) the pointwise penalized objective attains its minimum, and
its minimizing set has smallest and largest transport costs \(D_0^-\) and \(D_0^+\).

The source gives the following necessary and sufficient alternatives for primal attainment:
\begin{enumerate}[leftmargin=2em,label=(\alph*)]
\item some dual minimizer satisfies \(\lambda^*>\kappa\);
\item \(\lambda^*=\kappa>0\) is the unique dual minimizer, pointwise minimizers exist almost
      surely, and
\[
 \int D_0^-(\kappa,\zeta)d\nu
 \le\theta^p\le
 \int D_0^+(\kappa,\zeta)d\nu;
\]
\item \(\lambda^*=\kappa=0\) is the unique minimizer,
      \(\operatorname*{argmax}\Psi\ne\varnothing\), and
\[
 \int D_0^-(0,\zeta)d\nu\le\theta^p.
\]
\end{enumerate}
When one of these holds, measurable exact selectors \(T^+,T^-\) and a scalar
\(q\in[0,1]\) give
\begin{equation}\label{eq:gk-worst-law}
 \mu^*=q(T^+)_\#\nu+(1-q)(T^-)_\#\nu,
\end{equation}
with the cost interpolated to the budget.  The associated coupling is supported on
\[
 \{(\xi,\zeta):
   \xi\in\operatorname*{argmin}_{u}
      [\lambda^*d^p(u,\zeta)-\Psi(u)]\}.
\]
If \(\lambda^*>0\), complementary slackness forces the transport budget to be saturated and this
coupling is optimal in the definition of \(W_p\).

Conversely, if a worst-case law exists, equality in weak duality implies both complementary
slackness and concentration on the pointwise argmin relation.  Disintegrating an optimal coupling
and using the extreme costs of its conditional support yields the budget inequalities above.  This
proves necessity.

If \(\Xi\) is convex, \(\Psi\) is concave, and \(d^p(\cdot,\zeta)\) is convex almost surely, each
pointwise argmin set is convex.  The conditional barycenter of an optimal coupling then remains in
the argmin set by Jensen's inequality.  Thus randomization can be removed: a single measurable
map \(T^*\) satisfies \((T^*)_\#\nu\) primal optimal.

\section{Corollary 2: empirical nominal distributions}
Let
\[
 \nu=\frac1N\sum_{i=1}^N\delta_{\widehat\xi_i}.
\]
The dual formula becomes
\[
 v_D=\inf_{\lambda\ge0}
 \left\{\lambda\theta^p+
 \frac1N\sum_{i=1}^N
   \sup_{\xi\in\Xi}
    [\Psi(\xi)-\lambda d^p(\xi,\widehat\xi_i)]\right\}.
\]
The constructive proof represents an \(\varepsilon\)-optimal or optimal law as
\[
 \mu=\frac1N\sum_{i=1}^N\sum_{k=1}^{K_i}p_{ik}\delta_{\xi_{ik}},
 \qquad \sum_kp_{ik}=1,
\]
where each \(\xi_{ik}\) is a pointwise near optimizer for the \(i\)-th nominal atom.  The remaining
choice of weights is a finite linear program with one global transport-budget constraint.  At an
extreme point, all but at most one nominal atom use a single destination; hence an optimal law, when
it exists, can be represented using at most \(N+1\) effective atoms.  Without attainment, the same
construction gives \(\varepsilon\)-optimal laws with at most \(2N\) displayed atoms.

\section{Source-faithfulness remarks}
\begin{enumerate}[leftmargin=2em]
\item The paper's proof is not an invocation of a general Kantorovich duality theorem.  It is a
constructive scalar-dual argument based on measurable near minimizers and interpolation of their
transport costs.
\item The boundary case \(\lambda^*=\kappa\) is essential.  Replacing it by an interior Slater
argument would lose one of the main contributions of the source.
\item The selectors are measurable only after passing to the \(\nu\)-completion, exactly as in the
source.
\item The notation for upper and lower versions of \(D,D_0\) is typographically difficult in the
published PDF.  Superscripts \(+\) and \(-\) are used here to make the ordering visible without
changing the definitions.
\end{enumerate}

\begin{thebibliography}{99}
\bibitem{ChenGeorgiouPavon2021} Yongxin Chen, Tryphon T. Georgiou, and Michele Pavon. \emph{Stochastic Control Liaisons: Richard Sinkhorn Meets Gaspard Monge on a Schr\"odinger Bridge}. SIAM Review, 63(2):249--313, 2021. doi:10.1137/20M1339982. arXiv:2005.10963.
\bibitem{Leonard2014Survey} Christian L\'eonard. \emph{A Survey of the Schr\"odinger Problem and Some of Its Connections with Optimal Transport}. Discrete and Continuous Dynamical Systems--A, 34(4):1533--1574, 2014. doi:10.3934/dcds.2014.34.1533. arXiv:1308.0215.
\bibitem{Fortet1940} Robert Fortet. \emph{Resolution d'un systeme d'equations de M. Schrodinger}. Journal de mathematiques pures et appliquees, 9e serie, 19(1--4):83--105, 1940.
\bibitem{SinkhornKnopp1967} Richard Sinkhorn and Paul Knopp. \emph{Concerning Nonnegative Matrices and Doubly Stochastic Matrices}. Pacific Journal of Mathematics, 21(2):343--348, 1967.
\bibitem{FranklinLorenz1989} Joel Franklin and Jens Lorenz. \emph{On the Scaling of Multidimensional Matrices}. Linear Algebra and its Applications, 114/115:717--735, 1989. doi:10.1016/0024-3795(89)90490-4.
\bibitem{Carroll2004} Joseph E. Carroll. \emph{Birkhoff's Contraction Coefficient}. Linear Algebra and its Applications, 389:227--234, 2004. doi:10.1016/j.laa.2004.02.039.
\bibitem{EvesonNussbaum1995} Simon P. Eveson and Roger D. Nussbaum. \emph{An Elementary Proof of the Birkhoff--Hopf Theorem}. Mathematical Proceedings of the Cambridge Philosophical Society, 117(1):31--55, 1995.
\bibitem{BlanchetMurthy2019} Jose Blanchet and Karthyek R. A. Murthy. \emph{Quantifying Distributional Model Risk via Optimal Transport}. Mathematics of Operations Research, 44(2):565--600, 2019. doi:10.1287/moor.2018.0936. arXiv:1604.01446.
\bibitem{GaoKleywegt2023} Rui Gao and Anton J. Kleywegt. \emph{Distributionally Robust Stochastic Optimization with Wasserstein Distance}. Mathematics of Operations Research, 48(2):603--655, 2023. doi:10.1287/moor.2022.1275. arXiv:1604.02199.
\bibitem{MohajerinEsfahaniKuhn2018} Peyman Mohajerin Esfahani and Daniel Kuhn. \emph{Data-Driven Distributionally Robust Optimization Using the Wasserstein Metric: Performance Guarantees and Tractable Reformulations}. Mathematical Programming, 171(1--2):115--166, 2018. doi:10.1007/s10107-017-1172-1. arXiv:1505.05116.
\bibitem{WangGaoXie2025} Jie Wang, Rui Gao, and Yao Xie. \emph{Sinkhorn Distributionally Robust Optimization}. Operations Research, 2025. doi:10.1287/opre.2023.0294. arXiv:2109.11926.
\bibitem{Girsanov1960} I. V. Girsanov. \emph{On Transforming a Certain Class of Stochastic Processes by Absolutely Continuous Substitution of Measures}. Theory of Probability and Its Applications, 5(3):285--301, 1960. doi:10.1137/1105027.
\bibitem{CameronMartin1945} R. H. Cameron and W. T. Martin. \emph{Transformations of Wiener Integrals under a General Class of Linear Transformations}. Transactions of the American Mathematical Society, 58(2):184--219, 1945.
\bibitem{Yeh1978} J. Yeh. \emph{Transformation of Conditional Wiener Integrals under Translation and the Cameron--Martin Translation Theorem}. Tohoku Mathematical Journal, 30(4):505--515, 1978.
\bibitem{Kakutani1948} Shizuo Kakutani. \emph{On Equivalence of Infinite Product Measures}. Annals of Mathematics, 49(1):214--224, 1948.
\bibitem{ShiDeBortoliCampbellDoucet2023} Yuyang Shi, Valentin De Bortoli, Andrew Campbell, and Arnaud Doucet. \emph{Diffusion Schr\"odinger Bridge Matching}. Advances in Neural Information Processing Systems, 2023. arXiv:2303.16852.
\bibitem{LiuLipmanNickelKarrerTheodorouChen2024} Guan-Horng Liu, Yaron Lipman, Maximilian Nickel, Brian Karrer, Evangelos A. Theodorou, and Ricky T. Q. Chen. \emph{Generalized Schr\"odinger Bridge Matching}. International Conference on Learning Representations, 2024. arXiv:2310.02233.
\bibitem{DomingoEnrichDrozdzalKarrerChen2025} Carles Domingo-Enrich, Michal Drozdzal, Brian Karrer, and Ricky T. Q. Chen. \emph{Adjoint Matching: Fine-Tuning Flow and Diffusion Generative Models with Memoryless Stochastic Optimal Control}. arXiv preprint arXiv:2409.08861, 2025.

\bibitem{Birkhoff1957} Garrett Birkhoff. \emph{Extensions of Jentzsch's Theorem}. Transactions of the American Mathematical Society, 85(1):219--227, 1957.
\bibitem{Sinkhorn1964} Richard Sinkhorn. \emph{A Relationship Between Arbitrary Positive Matrices and Doubly Stochastic Matrices}. Annals of Mathematical Statistics, 35(2):876--879, 1964.
\bibitem{Sinkhorn1967Prescribed} Richard Sinkhorn. \emph{Diagonal Equivalence to Matrices with Prescribed Row and Column Sums}. American Mathematical Monthly, 74(4):402--405, 1967.
\bibitem{Bauer1965} Friedrich L. Bauer. \emph{An Elementary Proof of the Hopf Inequality for Positive Operators}. Numerische Mathematik, 7:331--337, 1965.
\bibitem{Bushell1973} P. J. Bushell. \emph{Hilbert's Metric and Positive Contraction Mappings in a Banach Space}. Archive for Rational Mechanics and Analysis, 52:330--338, 1973.

\end{thebibliography}

\end{document}
